\documentclass[11pt]{article}
\usepackage[a4paper,margin=28mm]{geometry}
\usepackage{amsmath,amssymb,amsthm}
\usepackage[colorlinks=true,linkcolor=blue,urlcolor=blue,citecolor=blue]{hyperref}
\newtheorem*{theorem*}{Theorem}
\newtheorem*{corollary*}{Corollary}
\newenvironment{statement}[1]{\par\medskip\noindent\phantomsection\textbf{#1}\ }{\par\medskip}
\setlength{\emergencystretch}{3em}
\newcommand{\Q}{\mathbb Q}
\newcommand{\Z}{\mathbb Z}
\newcommand{\F}{\mathbb F}
\DeclareMathOperator{\End}{End}
\DeclareMathOperator{\Gal}{Gal}
\DeclareMathOperator{\GL}{GL}
\DeclareMathOperator{\PGL}{PGL}
\DeclareMathOperator{\Hom}{Hom}
\DeclareMathOperator{\tr}{tr}
\title{Rational \(7\)-torsion and the \(2\)-rank of geometrically split abelian surfaces}
\author{}
\date{10 November 2026}

\begin{document}
\maketitle

\begin{theorem*}
Let \(A/\Q\) be an abelian surface that is simple over \(\Q\) and geometrically split. If
\(A(\Q)\) has a point of order \(7\), then \(A(\Q)[2]\cong(\Z/2)^r\) with \(r\le2\).
\end{theorem*}

\begin{corollary*}
No genus-\(2\) Jacobian over \(\Q\) that is geometrically split and simple over \(\Q\) has a
rational torsion subgroup containing \((\Z/2)^3\times\Z/7\). In particular,
\href{https://genus-2-torsion.github.io/group.html?g=2.2.14}{\([2,2,14]\)} does not occur as the
rational torsion subgroup of such a Jacobian. Together with the
\href{https://genus-2-torsion.github.io/proofs/qsplit-two-rank-obstruction.html}{\(2\)-rank
obstruction for surfaces split over \(\Q\)}, no geometrically split genus-\(2\) Jacobian over
\(\Q\) has rational torsion subgroup \([2,2,14]\).
\end{corollary*}

Here ``geometrically split'' means that \(A_{\bar\Q}\) is isogenous to a product of two elliptic
curves. We impose no restriction on the degree of this isogeny, on the field of definition of the
geometric endomorphisms or on the integral endomorphism ring, and we do not assume that \(A\) is
principally polarized.

The proof distinguishes the cases of the
\href{https://genus-2-torsion.github.io/proofs/qsimple-63-obstruction.html}{proof for order 63},
whose results we use (Section~\ref{inputs}). If the elliptic curves in \(A_{\bar\Q}\) are defined
over a quadratic field, the \(2\)-primary torsion of \(A(\Q)\) is controlled by that of two
isogenous elliptic curves over this field, one of which has a point of order \(7\); Kenku and
Momose's theorem then bounds it (Section~\ref{quadratic}). The CM case follows from a point count
at \(3\) (Section~\ref{cm}). By the proof for order 63, the only remaining endomorphism fields are
cyclic of degree \(3\) or \(6\). There the Galois action on the Tate module has a determinant over
\(\Q(\sqrt{-3})\) whose reduction modulo \(2\) has order \(3\), and this is incompatible with
\(2\)-rank \(3\) (Section~\ref{cyclic}).

\section{Conventions and inputs}\label{inputs}

All number fields are subfields of a fixed algebraic closure \(\bar\Q\). For a number field \(k\)
we write \(G_k=\Gal(\bar\Q/k)\). For a prime \(\ell\), \(\chi_\ell\) is the \(\ell\)-adic
cyclotomic character, \(T_\ell\) is the Tate module and \(V_\ell=T_\ell\otimes\Q_\ell\). For an
abelian variety \(X\), \(\End^0=\End\otimes\Q\) and \(\Hom^0=\Hom\otimes\Q\). For a variety \(Y\)
over \(\bar\Q\) and \(g\in G_{\Q}\), \(Y^g\) is the conjugate variety, \(g:Y(\bar\Q)\to
Y^g(\bar\Q)\) the induced bijection, and a morphism \(f:Y\to Z\) has the conjugate
\(g(f):Y^g\to Z^g\), characterized by \(g(f)(g(P))=g(f(P))\). For \(K\) a quadratic field we write
\(\sigma\) for its nontrivial automorphism.

We use the following.
\begin{itemize}
\item \textbf{(I1) Kenku--Momose} \cite[Theorem~(0.1)]{KM}: let \(N\) be a product of powers of
\(2,3,5,7,11,13\), and let \(m\mid N\). If \(X_1(m,N)\) has positive genus and \((m,N)\) is not one
of \((1,11)\), \((1,13)\), \((1,14)\), \((1,15)\), \((1,16)\), \((1,18)\), \((2,10)\), \((2,12)\),
then \(Y_1(m,N)(k)=\emptyset\) for every quadratic field \(k\). The curves \(X_1(28)=X_1(1,28)\) and
\(X_1(2,14)\) have genus \(10\) and \(4\). Hence \emph{for every quadratic field \(k\) and every elliptic curve
\(E\) over \(k\), the group \(E(k)\) has no subgroup isomorphic to \(\Z/28\) or to
\(\Z/2\times\Z/14\).}

\item \textbf{(I2) Serre--Tate}, as stated in Section~1 of the
\href{https://genus-2-torsion.github.io/proofs/qsimple-63-obstruction.html#inputs}{proof for
order 63}: let \(K\) be a finite extension of \(\Q_p\) with residue field \(k\), and \(A\) an
abelian variety over \(K\). (a) Good reduction is detected by the inertia action on
\(T_\ell(A)\), \(\ell\neq p\), and is preserved by \(K\)-isogenies \cite[\S1, Theorem~1 and
Corollary~2]{ST}. (b) If \(A\) has good reduction \(\widetilde A\) and \(m\) is prime to \(p\),
then \(A(K)[m]\) injects into \(\widetilde A(k)\) \cite[Lemma~2]{ST}. (c) If \(A\) has
potentially good reduction, then it has good reduction over a finite extension of \(K\) with
residue field \(k\) \cite[\S2]{ST}. For an abelian surface \(B\) over \(\F_q\),
\(\#B(\F_q)\le(1+\sqrt q)^4\) by Weil's theorem.

\item \textbf{(I3) The case distinction.} Let \(A/\Q\) be a geometrically split abelian surface.
Exactly one of the following holds: \textbf{(N)} \(A_{\bar\Q}\) is isogenous to \(E_1\times E_2\)
with \(E_1\), \(E_2\) not isogenous; \textbf{(CM)} \(A_{\bar\Q}\) is isogenous to \(E^2\) with
\(E\) a CM curve; \textbf{(S)} \(A_{\bar\Q}\) is isogenous to \(E^2\) with \(E\) without CM. In case
(S), \(D=\End^0(A_{\bar\Q})\cong M_2(\Q)\); \(L\) is the smallest field over which all geometric
endomorphisms of \(A\) are defined, a finite Galois extension of \(\Q\); and
\(\Gamma=\Gal(L/\Q)\) acts faithfully on \(D\), the element \(g\) acting by
\(\varphi\mapsto g(\varphi)=R_g\varphi R_g^{-1}\) with \(R_g\in D^\times\) unique up to
\(\Q^\times\). Thus \(\Gamma\) embeds in \(\PGL_2(\Q)\), and \(\End^0_{\Q}(A)=D^\Gamma\). An element
\(r\in\Gamma\) of order \(n\in\{3,4,6\}\) has a unique representative \(R\) with \(\tr R=1\), and
\(R^2-R+1/c=0\) with \(c=1,2,3\) respectively. (Sections~2 and~5 of the
\href{https://genus-2-torsion.github.io/proofs/qsimple-63-obstruction.html}{proof for order 63}.)

\item \textbf{(I4)} \cite[Lemma~2]{P63} Let \(X\) be an abelian surface over a number field
\(k\), \(E\subset X\) an elliptic curve defined over \(k\), and \(Q\in X(k)\) of prime order
\(q\). Then either \(Q\in E(k)\), or the image of \(Q\) in \((X/E)(k)\) has order \(q\).

\item \textbf{(I5)} \cite[Lemma~3]{P63} In case (S), every elliptic curve in \(A_{\bar\Q}\) is
defined over \(L\), and any two of them are \(L\)-isogenous.

\item \textbf{(I6)} \cite[Proposition~8]{P63} Let \(A\) be of type (S) and simple over \(\Q\),
and suppose that \(A(\Q)\) has a point of order \(7\). Then \(\Gamma\) is cyclic of order \(2\),
\(3\) or \(6\). The standing hypothesis of that proof is a point of order \(63\), but the proof of
Proposition~8 (its Section~7, together with \(\Gamma\neq C_1\) from Section~5) uses only simplicity
over \(\Q\) and the point of order \(7\), through Lemmas~6 and~7 there with \(p=7\) and their common input, Lemma~5, which uses a
polarization over \(\Q\) of arbitrary degree.

\item \textbf{(I7) Lattices and pairings.} For a continuous representation of a profinite group
on a finite-dimensional \(\Q_\ell\)-vector space, the composition factors of the reduction of a
stable lattice do not depend on the lattice (Brauer--Nesbitt); the same holds for two isomorphic
representations, after transporting a lattice through the isomorphism. For an isogeny \(u:E\to E'\) of
elliptic curves and the Weil pairings \(e\) on \(V_\ell\), \(e(ux,uy)=\deg(u)\,e(x,y)\); this extends to
quasi-isogenies, with \(\deg(u/m)=\deg(u)/m^2\). Moreover
\(e(gx,gy)=\chi_\ell(g)\,e(x,y)\) for \(g\in G_{\Q}\), where \(g:V_\ell(E)\to V_\ell(E^g)\).
\end{itemize}

\section{Elliptic curves over a quadratic field}\label{quadratic}

\begin{statement}{Lemma 1.}\label{lemma-km}
Let \(K\) be a quadratic field and let \(E'\) and \(E''\) be \(K\)-isogenous elliptic curves over
\(K\). If \(E'(K)\) has a point of order \(7\), then \(E''(K)[2^\infty]\) has order at most \(2\).
\end{statement}

\emph{Proof.} The \(G_K\)-module \(E'[7]\) has the trivial character as a composition factor, and
since its determinant is the cyclotomic character \(\bar\chi_7\), its composition factors are
\(\mathbf 1\) and \(\bar\chi_7|_{G_K}\). By (I7), applied to the isomorphism \(V_7(E')\cong V_7(E'')\) induced by the \(K\)-isogeny, so
are those of \(E''[7]\). Let
\(C\subset E''[7]\) be a \(G_K\)-stable line. If \(G_K\) acts trivially on \(C\), put
\(E'''=E''\). Otherwise \(G_K\) acts on \(C\) by \(\bar\chi_7\), hence trivially on
\(E''[7]/C\); put \(E'''=E''/C\); it contains the subgroup \(E''[7]/C\) of order \(7\), whose points
are \(K\)-rational because \(G_K\) acts trivially on it. In both cases \(E'''(K)\) has a point of order \(7\), and the \(K\)-isogeny
\(E''\to E'''\) has degree \(1\) or \(7\), so it maps \(E''[2^\infty]\) isomorphically and
\(G_K\)-equivariantly onto \(E'''[2^\infty]\). If \(E''(K)[2^\infty]\) had order at least \(4\), it
would contain \(\Z/4\) or \(\Z/2\times\Z/2\), and \(E'''(K)\) would contain a subgroup isomorphic to
\(\Z/28\) or to \(\Z/2\times\Z/14\), contradicting (I1). \qed

\begin{statement}{Proposition 2.}\label{prop-quadratic}
Let \(A\) be an abelian surface over \(\Q\), \(K\) a quadratic field and \(E\subset A_K\) an
elliptic curve defined over \(K\) with \(E^\sigma\neq E\). If \(A(\Q)\) has a point of order
\(7\), then \(A(\Q)[2^\infty]\) has order at most \(4\).
\end{statement}

\emph{Proof.}
\begin{enumerate}
\item \emph{A sequence.} Choose \(\tilde\sigma\in G_{\Q}\) restricting to \(\sigma\) on \(K\). The
curve \(E^\sigma=\tilde\sigma(E)\subset A_K\) is defined over \(K\) and different from \(E\), so
\(E\cap E^\sigma\) is finite. Let \(\pi:A_K\to Y:=A_K/E^\sigma\) be the quotient map; it is
defined over \(K\), and its restriction to \(E\) has the finite kernel \(E\cap E^\sigma\), so it is
a \(K\)-isogeny \(E\to Y\). If \(P\in A(\Q)\) and \(\pi(P)=0\), then \(P\in E^\sigma\), hence
\(P=\tilde\sigma(P)\in\tilde\sigma(E^\sigma)=E\), and \(P\in E(K)\). So the sequence
\[0\to A(\Q)[2^\infty]\cap E(K)\to A(\Q)[2^\infty]\xrightarrow{\ \pi\ }Y(K)[2^\infty]\]
is exact, and \(\#A(\Q)[2^\infty]\le\#E(K)[2^\infty]\cdot\#Y(K)[2^\infty]\).

\item \emph{A point of order \(7\).} Let \(P_7\in A(\Q)\) have order \(7\). By (I4) over \(K\),
applied to \(E^\sigma\subset A_K\), either \(P_7\in E^\sigma\), and then \(P_7\in E(K)\) by step 1,
or \(\pi(P_7)\in Y(K)\) has order \(7\). So \(E(K)\) or \(Y(K)\) has a point of order \(7\).

\item Since \(E\) and \(Y\) are \(K\)-isogenous, \hyperref[lemma-km]{Lemma~1} gives
\(\#E(K)[2^\infty]\le2\) and \(\#Y(K)[2^\infty]\le2\), and step 1 gives
\(\#A(\Q)[2^\infty]\le4\). \qed
\end{enumerate}

\medskip
\textbf{Type (N).} Let \(A\) be of type (N) and simple over \(\Q\). Then \(A_{\bar\Q}\) contains
exactly two elliptic curves; \(G_{\Q}\) exchanges them, since otherwise they would be defined over
\(\Q\); so they are defined over a quadratic field \(K\), and \(\sigma\) exchanges them
(Section~3 of the
\href{https://genus-2-torsion.github.io/proofs/qsimple-63-obstruction.html#nonisogenous}{proof
for order 63}). \hyperref[prop-quadratic]{Proposition~2} applies.

\medskip
\textbf{Type (S) with \(\Gamma=C_2\).} Let \(A\) be of type (S), simple over \(\Q\), with
\(\Gamma\cong C_2\). Then \(K:=L\) is quadratic, and every elliptic curve \(E\subset A_{\bar\Q}\) is
defined over \(K\) by (I5). If \(E^\sigma=E\), then \(E\) would be stable under \(G_{\Q}\), hence
an elliptic curve in \(A\) defined over \(\Q\), and \(A\) would not be simple over \(\Q\). So
\(E^\sigma\neq E\), and \hyperref[prop-quadratic]{Proposition~2} applies.

\section{The CM case: a point count at \(3\)}\label{cm}

Let \(A\) be of type (CM), and let \(P_7\in A(\Q)\) have order \(7\). The CM curve \(E\) has
integral \(j\)-invariant and hence potentially good reduction at every prime \cite[\S2]{ST}. Over a finite
extension of \(\Q_3\), \(A\) becomes isogenous to \(E^2\) and \(E\) acquires good reduction; by
(I2)(a), \(A\) acquires good reduction there. By (I2)(c), there is a finite extension \(K'\) of
\(\Q_3\) with residue field \(\F_3\) over which \(A\) has good reduction \(\widetilde A\). The
subgroup of \(A(\Q)\subseteq A(K')\) generated by \(P_7\) and \(A(\Q)[2^\infty]\) has order
\(7\cdot\#A(\Q)[2^\infty]\), which is prime to \(3\); by (I2)(b) it injects into
\(\widetilde A(\F_3)\). Hence
\[7\cdot\#A(\Q)[2^\infty]\ \le\ (1+\sqrt3)^4=28+16\sqrt3\ <\ 56,\]
because \(16\sqrt3<28\) is equivalent to \(768<784\). So \(\#A(\Q)[2^\infty]\le4\). This case does
not use simplicity over \(\Q\).

\section{Type (S) with \(\Gamma\) cyclic of order \(3\) or \(6\)}\label{cyclic}

In this section \(A\) is of type (S), \(\Gamma=\langle r\rangle\) is cyclic of order
\(n\in\{3,6\}\), and \(R\) is the trace-one representative of \(r\), with \(R^2-R+1/c=0\) and
\(c=1\) or \(3\). Put \(F=\Q[R]\subset D\). The discriminant of \(t^2-t+1/c\) is \(-3\) or
\(-1/3\), so \(F\) is a field isomorphic to \(\Q(\sqrt{-3})\). Since \(R\) is not a scalar, its
centralizer in \(D\) is \(F\), so \(\End^0_{\Q}(A)=D^\Gamma=F\). We write \(\bar x\) for the
conjugate of \(x\in F\) and \(N(x)=x\bar x\); thus \(\bar R=1-R\) and \(N(R)=\det R=1/c>0\). The
two roots \(R,\bar R\) of \(t^2-t+1/c\) are the eigenvalues of \(R\), so \(R/\bar R\) is a
primitive \(n\)-th root of unity (Section~5 of the
\href{https://genus-2-torsion.github.io/proofs/qsimple-63-obstruction.html#endomorphism-field}{proof
for order 63}); explicitly \(R/\bar R=(1\pm\sqrt{-3})/2\) when \(n=6\), and \((-1\pm\sqrt{-3})/2\)
when \(n=3\).

\begin{statement}{Lemma 3 (determinant).}\label{lemma-det}
Let \(\ell\) be a prime. Then \(V_\ell(A)\) is a free \(F\otimes\Q_\ell\)-module of rank \(2\), on
which \(G_{\Q}\) acts \(F\otimes\Q_\ell\)-linearly, and the determinant
\(\delta_\ell:G_{\Q}\to(F\otimes\Q_\ell)^\times\) of this action is \(\delta_\ell=\chi_\ell\,
\varepsilon\), where \(\varepsilon:G_{\Q}\to F^\times\) is the character that factors through
\(\Gamma\) and maps \(r\) to \(R/\bar R\). In particular \(\varepsilon\) has order \(n\).
\end{statement}

\emph{Proof.}
\begin{enumerate}
\item \emph{The multiplicity space.} Choose an elliptic curve \(E\subset A_L\); it is defined over
\(L\) by (I5). Put \(H=\Hom^0(E_{\bar\Q},A_{\bar\Q})\), a \(\Q\)-vector space of dimension \(2\).
Every element of \(H\) is defined over \(L\): if \(p\in\Hom^0(A_L,E)\) is a quasi-projection defined over
\(L\) (that is, \(p\circ i=1\) for the inclusion \(i\); for instance the map induced by a
rank-one idempotent \(e\in D\) with image \(E\), which exists by Poincar\'e reducibility, as the
projection along a complement of \(E\), and is defined over \(L\) like every element of \(D\)), then \(h\mapsto h\circ p\) embeds \(H\)
\(G_L\)-equivariantly into \(D\), on which \(G_L\) acts trivially. The algebra \(D\) acts
on \(H\) by composition, faithfully because \(D\) is simple and the action is unital; so \(D=\End_{\Q}(H)\) as \(\dim_{\Q}H=2\), and \(H\) is a
one-dimensional \(F\)-vector space. The evaluation map
\[\Phi:V_\ell(E)\otimes_{\Q}H\to V_\ell(A),\qquad v\otimes h\mapsto h(v),\]
is an isomorphism: for a basis \(h_1,h_2\) of \(H\), a nonzero integer multiple of \((h_1,h_2):E^2\to A_{\bar\Q}\)
is a homomorphism, and it is an isogeny (otherwise its image would be an elliptic curve \(E_1\), and
\(h_1,h_2\) would be linearly dependent in the one-dimensional space \(\Hom^0(E,E_1)\)); \(\Phi\) is
the map it induces, up to that integer. The
map \(\Phi\) is \(F\)-linear, \(F\) acting on \(H\). Hence \(V_\ell(A)\cong
V_\ell(E)\otimes_{\Q_\ell}(F\otimes\Q_\ell)\) is free of rank \(2\) over \(F\otimes\Q_\ell\). As
\(F=\End^0_{\Q}(A)\), the action of \(G_{\Q}\) is \(F\otimes\Q_\ell\)-linear.

\item \emph{The action of \(g\in G_{\Q}\).} The curve \(E^g\subset A_L\) is \(L\)-isogenous to
\(E\) by (I5); choose an invertible \(u\in\Hom^0(E,E^g)\). Define
\(B:V_\ell(E)\to V_\ell(E)\) by \(B(v)=u^{-1}(g(v))\) and \(C:H\to H\) by \(C(h)=g(h)\circ u\).
For \(v\in V_\ell(E)\) and \(h\in H\),
\[g\bigl(h(v)\bigr)=g(h)\bigl(g(v)\bigr)=g(h)\bigl(u(Bv)\bigr)=C(h)(Bv),\]
so \(g\circ\Phi=\Phi\circ(B\otimes C)\). For \(\varphi\in D\) we have
\(C(\varphi\circ h)=g(\varphi)\circ C(h)\); so \(C\in\GL(H)=D^\times\) satisfies
\(g(\varphi)=C\varphi C^{-1}\). In particular \(C\) commutes with \(F=D^\Gamma\), so \(C\) is
multiplication by an element \(\gamma\in F^\times\), and
\[\delta_\ell(g)=\det(B)\,\gamma^2.\]
By (I7), \(e(Bv,Bw)=\deg(u)^{-1}e(g(v),g(w))=\deg(u)^{-1}\chi_\ell(g)\,e(v,w)\) for the Weil
pairing \(e\) of \(E\), so \(\det B=\chi_\ell(g)/\deg(u)\). Hence
\[\delta_\ell(g)=\chi_\ell(g)\,\varepsilon(g),\qquad \varepsilon(g):=\gamma^2/\deg(u)\in F^\times.\]
The element \(\varepsilon(g)=\delta_\ell(g)/\chi_\ell(g)\) does not depend on the choice of \(u\),
and its definition does not involve \(\ell\).

\item \emph{\(\varepsilon\) is a character of \(\Gamma\).} It is a homomorphism, being a quotient of
two homomorphisms to the commutative group \((F\otimes\Q_\ell)^\times\). For \(g\in G_L\) we may
take \(u=1\), since \(E^g=E\); then \(C(h)=g(h)=h\) for all \(h\in H\), so \(\gamma=1\) and
\(\varepsilon(g)=1\). Hence \(\varepsilon\) factors through \(\Gamma\), and its values are roots of
unity in \(F\), of norm \(1\).

\item \emph{The value at \(r\).} Let \(g\) restrict to \(r\). Both \(C\) and \(R\) induce the
action of \(g\) on \(D\), so \(\gamma=\lambda R\) with \(\lambda\in\Q^\times\) (Skolem--Noether). Then
\(\varepsilon(g)=\mu R^2\) with \(\mu=\lambda^2/\deg(u)>0\), and
\(1=N(\varepsilon(g))=\mu^2N(R)^2\) gives \(\mu=1/N(R)\). So
\(\varepsilon(g)=R^2/(R\bar R)=R/\bar R\). \qed
\end{enumerate}

Step~4 uses \(N(R)=1/c>0\). For \(\Gamma=C_2\) the same computation, with \(R^2=d\) and
\(N(R)=-d\), gives \(\varepsilon=1\) when \(\End^0_{\Q}(A)=\Q(\sqrt d)\) is real; this is one
reason why \(\Gamma=C_2\) is treated by Proposition~2 instead.

\begin{statement}{Proposition 4.}\label{prop-cyclic}
Let \(A\) be of type (S) with \(\Gamma\) cyclic of order \(3\) or \(6\). Then the
\(\F_2[G_{\Q}]\)-module \(A[2]\) has a composition factor on which \(G_{\Q}\) acts nontrivially.
In particular \(A(\Q)[2]\) has rank at most \(2\).
\end{statement}

\emph{Proof.} Take \(\ell=2\). Since \(-3\equiv5\pmod8\), the prime \(2\) is inert in \(F\), so
\(F\otimes\Q_2\) is the unramified quadratic extension of \(\Q_2\); let \(\mathcal O\) be its ring
of integers, with residue field \(\F_4\). The lattice \(T'=\mathcal O\,T_2(A)\subset V_2(A)\) is
\(G_{\Q}\)-stable and is a free \(\mathcal O\)-module of rank \(2\). So \(M=T'/2T'\) is a
two-dimensional \(\F_4\)-vector space on which \(G_{\Q}\) acts \(\F_4\)-linearly, with
\(\F_4\)-determinant the reduction of \(\delta_2=\chi_2\varepsilon\) modulo \(2\mathcal O\). As
\(\chi_2\) takes values in \(1+2\Z_2\), this determinant is the reduction \(\bar\varepsilon\) of
\(\varepsilon\). By \hyperref[lemma-det]{Lemma~3}, \(\varepsilon(r)^2\) is a primitive cube root of
unity \(\zeta\), and \(\zeta\not\equiv1\pmod{2\mathcal O}\) because \(N(1-\zeta)=3\) is a unit at
\(2\). So \(\bar\varepsilon\neq1\).

Suppose that all composition factors of \(A[2]=T_2(A)/2T_2(A)\) are trivial. By (I7), applied to
the lattices \(T_2(A)\) and \(T'\), all composition factors of the \(\F_2[G_{\Q}]\)-module \(M\)
are trivial. In a basis adapted to a composition series, the image of \(G_{\Q}\) in
\(\GL_{\F_2}(M)\) consists of unipotent upper triangular matrices, so it is a \(2\)-group. Its
image under the \(\F_4\)-determinant is then a \(2\)-group in \(\F_4^\times\), which has order
\(3\); so \(\bar\varepsilon=1\), a contradiction.

If \(A(\Q)[2]\) had rank at least \(3\), then \(G_{\Q}\) would act trivially on a
three-dimensional subspace of \(A[2]\) and, since \(\GL_1(\F_2)=1\), on the one-dimensional
quotient; all composition factors of \(A[2]\) would be trivial. \qed

\medskip
\hyperref[prop-cyclic]{Proposition~4} uses neither a point of order \(7\) nor simplicity over
\(\Q\).

\section{Proof of the Theorem and consequence for the census}\label{census}

\emph{Proof of the Theorem.} Let \(A/\Q\) be simple over \(\Q\) and geometrically split, with a
rational point of order \(7\). By (I3), \(A\) is of type (N), (CM) or (S). In type (N),
\(\#A(\Q)[2^\infty]\le4\) by Section~\ref{quadratic}; in type (CM), by Section~\ref{cm}. In type
(S), \(\Gamma\) is cyclic of order \(2\), \(3\) or \(6\) by (I6). If \(\Gamma=C_2\), then
\(\#A(\Q)[2^\infty]\le4\) by Section~\ref{quadratic}; if \(\Gamma\) has order \(3\) or \(6\), then
\(A(\Q)[2]\) has rank at most \(2\) by \hyperref[prop-cyclic]{Proposition~4}. In all cases
\(A(\Q)[2]\) has rank at most \(2\). \qed

\medskip
\emph{Proof of the Corollary.} The Jacobian of a curve of genus \(2\) over \(\Q\) is an abelian
surface over \(\Q\), and \((\Z/2)^3\times\Z/7\) has a point of order \(7\) and \(2\)-rank \(3\).
For Jacobians split over \(\Q\), the
\href{https://genus-2-torsion.github.io/proofs/qsplit-two-rank-obstruction.html}{\(2\)-rank
obstruction} shows that a rational point of order \(7\) forces \(2\)-rank at most \(2\). \qed

\medskip
The group \href{https://genus-2-torsion.github.io/group.html?g=2.2.14}{\([2,2,14]\)} is impossible
in the class ``geometrically split, simple over \(\Q\)''. It was already impossible in the class
``split over \(\Q\)''. Hence the combined geometrically split cell of
\([2,2,14]\), whose grade was open (``\(\infty\)?'' on the census), becomes impossible
(\(\emptyset\)).

The argument places no restriction on the geometric splitting degree, the field of definition of
the geometric endomorphisms, or the integral endomorphism ring, and it does not use a principal
polarization.

\section{Remarks}\label{remarks}

\begin{enumerate}
\item \emph{The group \([2,28]\).} In types (N) and (CM), and in type (S) with \(\Gamma=C_2\), the
proof gives \(\#A(\Q)[2^\infty]\le4\) when \(A(\Q)\) has a point of order \(7\). So a
geometrically split abelian surface over \(\Q\), simple over \(\Q\), whose rational torsion
contains \(\Z/2\times\Z/28\) is of type (S) with \(\Gamma\) cyclic of order \(3\) or \(6\). These
two cases are not settled here, and the census cell of
\href{https://genus-2-torsion.github.io/group.html?g=2.28}{\([2,28]\)} for geometrically split
Jacobians that are simple over \(\Q\) stays open.

\item \emph{Modular surfaces.} If \(A\) is the abelian surface \(A_f\) attached to a weight-\(2\)
newform \(f\) whose coefficient field is \(\End^0_{\Q}(A)\), then \(\delta_\ell\) is the
cyclotomic character times the nebentypus of \(f\), so \(\varepsilon\) corresponds to the
nebentypus or to its inverse, depending on the identification of \(F\) with the coefficient field;
in particular both have order \(|\Gamma|\). For
example, \(J_1(13)\) and \(J_1(18)\) are the surfaces of the newforms 13.2.e.a and 18.2.c.a, with
nebentypus of order \(6\) and \(3\); their Sato--Tate groups in the LMFDB (curves 169.a.169.1 and
324.a.648.1) are \(E_6\) and \(E_3\), whose component groups \(\Gamma\) are cyclic of order \(6\)
and \(3\). This is consistent with \hyperref[lemma-det]{Lemma~3}, and is not used in the proof.

\item \emph{The point of order \(7\) is needed.} Without it, rational \(2\)-rank \(3\) occurs on
\(\Q\)-simple, geometrically split Jacobians of type (N), of type (CM), and of type (S) with
\(\Gamma=C_2\): the LMFDB curves 2560.a.819200.1, 4096.b.65536.1 and 25600.f.512000.1 have rational
torsion subgroup \(\Z/2\times\Z/2\times\Z/2\) and Sato--Tate groups \(N(\mathrm{SU}(2)\times
\mathrm{SU}(2))\), \(J(C_2)\) and \(J(E_1)\). By \hyperref[prop-cyclic]{Proposition~4}, this cannot
happen when \(\Gamma\) is cyclic of order \(3\) or \(6\).
\end{enumerate}

\begin{thebibliography}{P63}

\bibitem[KM]{KM}
M. A. Kenku and F. Momose, \emph{Torsion points on elliptic curves defined over quadratic fields},
Nagoya Math. J. \textbf{109} (1988), 125--149, Theorem~(0.1).
\href{https://doi.org/10.1017/S0027763000002816}{doi:10.1017/S0027763000002816}.

\bibitem[P63]{P63}
Census proof of 21 October 2026, \emph{An obstruction to rational \(63\)-torsion on geometrically
split abelian surfaces},
\url{https://genus-2-torsion.github.io/proofs/qsimple-63-obstruction.html}: Lemmas~2, 3
and~5--7, Proposition~8, and Sections~1, 2, 3, 5, 6 and~7.

\bibitem[ST]{ST}
J.-P. Serre and J. Tate, \emph{Good reduction of abelian varieties}, Ann. of Math. (2)
\textbf{88} (1968), 492--517: Theorem~1 and Corollary~2 (p.~493), Lemma~2 (p.~495), \S2
(pp.~496 and~498). \href{https://doi.org/10.2307/1970722}{doi:10.2307/1970722}.

\end{thebibliography}

\end{document}
